Una balanza basada en un resorte con amortiguamiento viscoso tiene constantes desconocidas \(k>0\) y \(c\ge0\), iguales en todos los ensayos. Para masas patrón de \(1\,\mathrm{kg}\) y \(2\,\mathrm{kg}\) se observan frecuencias de oscilación de \(1.1\,\mathrm{Hz}\) y \(0.8\,\mathrm{Hz}\), respectivamente. Tratar estas medidas como datos exactos del modelo lineal.

\begin{enumerate}[label=(\alph*),ref=\alph*,leftmargin=*]
\item Determinar \(k\) y \(c\ge0\).
\item Obtener una fórmula para la masa a partir de una frecuencia \(f>0\), indicando cuándo hay cero, una o dos masas posibles.
\item Se mide \(f=0.2\,\mathrm{Hz}\) para un objeto cuya masa es mayor que \(1\,\mathrm{kg}\). Determinar su masa.
\end{enumerate}
